\documentclass[titlepage,a4paper]{jsarticle}
\usepackage{../../../sty/import}% 各種パッケージインポート
\usepackage{../../../sty/title}% タイトルページの変更
\usepackage{../../../sty/answergrid}
%% タイトルページの変数
% レポートタイトル
\title{心理学特論第7回目出席票}
% 提出日
\expdate{\today}
% 科目名
\subject{心理学特論}
% 分野
\class{情報経営システム工学分野}
% 学年
\grade{M1}
% 学籍番号
\mynumber{24336488}
% 記述者
\author{本間三暉}

\begin{document}
% titleページ作成
\maketitle
\section*{keyword\#1}
共有地の悲劇
\section*{keyword\#2}
ディレールメント
\section*{心理学特論第07回で取り上げた「社会的手抜き」と「比較優位」について考えてみましょう。
これまでの経験を振り返って簡潔にまとめて下さい。}
\section{「社会的手抜き」は、集団でやったことで手抜きが生じる現象です。集団において自分が手を抜いた、または手を抜かれたように感じたことがあれば、なるべく具体的に説明して下さい。}
私が社会的手抜きを感じた経験は，B3の授業で行ったグループ学習である．
グループで課題に取り組む際，人数が多くなるほど，「誰かがやってくれるだろう」という雰囲気が生まれやすく，一部の人だけが積極的に作業を進め，その他の人は受け身になる場面があった．
特に，資料の細かな修正や最終確認など，成果として見えにくい作業ほど担当が曖昧になりやすかった．

また，自分自身も，「まだ他の人が進めていないから後でやろう」と考えてしまい，主体的に動くのが遅れたことがある．
その結果，作業が締切直前に集中し，特定のメンバーに負担が偏ってしまった．
授業資料でも，社会的手抜きとは，集団で作業すると個人のときより無意識に力を抜いてしまう現象であると説明されており，自分の経験にも当てはまると感じた．

この経験から，集団作業では役割分担を明確にすることや，進捗共有を定期的に行うことが重要であると学んだ．
また，コミュニケーション不足が負担の偏りにつながることも実感した．

\section{自分にとっての比較優位を考えるために、スライド資料の「コンピテンシー」の部分を読んで、自分に該当しそうな要素について、なるべく具体的に、それこそ就活や進学でアピールするつもりで書いてみて下さい。}
私にとっての比較優位は，IT技術を活用して問題解決を行う力と，全体を見ながら改善点を考えられる点である．
授業では，比較優位とは単純な能力差ではなく，「それぞれが相対的に得意な分野を担当することで，全体の生産性を高める考え方」であると説明されていた．

私は，企業で技術系のアルバイトとしてWebシステムの開発に携わった経験があり，実際の利用を想定しながら機能改善や不具合修正を行ってきた．
また，学園祭実行委員会では，サーバやネットワークなどの運用に関わり，多くの人が利用するシステムを安定して動かすための作業を担当した．
さらに，個人でもWebアプリケーション開発を行っている．

これらの経験を通して，単に作業を進めるだけではなく，「利用者にとって使いやすいか」「運用時に問題が起きにくいか」を考えながら改善することを意識してきた．
また，問題が発生した際には，原因を整理しながら一つずつ切り分けて対応することを得意としている．

授業資料にあるコンピテンシーの中では，「イニシアチブ」「分析力」「自主独立性」「ストレス耐性」が自分に当てはまると考えている．
新しい技術について自分で調べながら試行錯誤し，問題が起きても粘り強く対応できる点が，自分の強みである．

\section{心理学特論課題サイコロワーク解答記入欄}
\AnswerGridFilled{
  {ウ}
  {イ}
  {ア}
  {ア}
  {ア}
  {エ}
  {イ}
  {ア}
  {エ}
  {イ}
  {比較優位}
  {絶対優位}
  {共有地の悲劇}
  {機会費用}
  {資料を作成する}
  {初頭効果}
  {ハイダー}
  {フェスティンガー}
  {原因帰属}
  {ANOVA}
}

\end{document}